% sections/07_discussion.tex

\section{Discussion}
\label{sec:discussion}

\subsection{Structure Over Scale}

The central result of this work is that reliability in agentic scientific ML
is achievable through architectural constraints rather than through model
capability alone.
The Agentic Chemometrician operates correctly---detecting leakage, refusing
to hallucinate metrics, routing non-routine decisions to a human---regardless
of which language model is used to drive the conversation.
This is possible because every tool in the system takes and returns typed,
schema-validated dataclasses.
A model cannot silently pass a leakage-inflated score downstream because the
\texttt{validate\_results} tool will serialize the group-structure mismatch
as a \texttt{ValidationWarning} in the contract, and that warning is a
first-class field that downstream tools inspect.
No prompt engineering is required to make this work; the enforcement lives in
the Python dataclass layer, not in the language model's inference.

The implication for the field is direct: the path to trustworthy LLM-assisted
scientific analysis runs through typed contracts and deterministic core logic,
not through next-generation foundation models.
A GPT-4-class model paired with an ungrounded tool API will repeat the same
failure modes we document in the introduction.
The same GPT-4-class model paired with the contract layer described here will
not---because the contract layer does not ask the model to check for leakage;
it checks automatically and surfaces the result as structured data.

\subsection{Modality Generality via the Declarative Registry}

A practical concern for any spectroscopic tool system is the proliferation of
instrument-specific branches.
NIR, FTIR, and Raman spectra differ in axis units, meaningful value ranges,
and which preprocessing operations are physically meaningful.
A naive implementation handles these differences with a cascade of
\texttt{if modality == "FTIR"} conditionals scattered across preprocessing,
validation, and reporting code.

The declarative \texttt{ModalityProfile} registry eliminates this pattern
by centralizing all modality-specific knowledge in a single frozen dataclass
per modality.
Adding support for a new modality---Raman variants, UV--Vis, or process mass
spectrometry---requires adding one \texttt{ModalityProfile} entry.
All downstream modules (preprocessing scope enforcement, axis-unit
validation, value-domain clipping) inherit the new behavior automatically.
The reliability mechanisms transfer with it: leakage detection, HITL gates,
and typed contracts do not need to be re-implemented or re-tested per modality.

This architectural property is especially important in industrial deployments
where the same analysis pipeline must service multiple product lines measured
by different instruments.

\subsection{Model Agnosticism}

The system is deliberately model-agnostic.
The MCP protocol does not specify which language model calls the tools, and
the tool implementations carry no assumptions about the model's capabilities
or output format.
This means the HITL gates and leakage-detection mechanism function the same
way whether the driving model is a frontier commercial LLM, an open-weight
model, or a rule-based dispatcher.
Reliability is not a property of the model; it is a property of the interface.

\subsection{When the HITL Cost Is Justified}

Human-in-the-loop approval adds latency and requires an expert to be
available at decision points~\citep{monarch2021hitl}.
This cost is justified under a specific set of conditions that match the
chemometrics domain closely:

\begin{itemize}
  \item \textbf{Low $n$, high stakes.} When sample counts are in the tens
    rather than thousands, a single leakage-inflated result can propagate
    through an entire method-development cycle and into a regulatory
    submission.
    The cost of one human review of a flagged validation result is negligible
    relative to the cost of discovering the error during an audit.
  \item \textbf{Leakage-prone data structure.} Replicate nesting---multiple
    scans per physical sample---is the default in spectroscopic protocols, not
    an edge case.
    Automatic leakage detection triggers human review precisely when it is
    most needed.
  \item \textbf{Model selection under ambiguity.} When the best-performing
    model is associated with validation warnings (small per-class counts,
    class imbalance, near-singular covariance), an autonomous system would
    either pick it without comment or suppress it with no recourse.
    A HITL gate surfaces both the metric and the warning and lets the
    domain expert decide, preserving the audit trail.
  \item \textbf{Novel datasets.} Method memory tools record approved
    preprocessing and model configurations across sessions, so HITL overhead
    decreases as institutional knowledge accumulates.
    The first analysis of a new product line incurs full human review;
    subsequent analyses of structurally similar datasets benefit from
    remembered, approved priors.
\end{itemize}

These conditions do not generalize to all ML settings.
For large-sample, low-stakes, well-characterized tasks, a fully autonomous
pipeline may be appropriate.
The system architecture is a deliberate choice for the domain, not a
claim that HITL is universally preferable.
